\section{Выводы}

\begin{enumerate}
    \item Построен воспроизводимый pipeline анализа loss landscape для CIFAR-10.
    \item Подтверждена гипотеза: в локальной окрестности решения полная и диагональная квадратики описывают поверхность лучше, чем FD-Hessian аппроксимация.
    \item Получены интерпретируемые связи между теорией (квадратичная локальная модель) и экспериментом (ошибки аппроксимации).
\end{enumerate}

\subsection{Ограничения}
\begin{itemize}
    \item Высокая стоимость полного прогона на CPU.
    \item Численная чувствительность FD-модели.
    \item Неполный набор сравнений по model2 в текущей версии отчёта.
\end{itemize}

\subsection{Следующие шаги}
\begin{itemize}
    \item Запустить системный sweep по $\epsilon$ для $q_3$.
    \item Добавить multi-seed анализ плоскостей $(d_1, d_2)$.
    \item Зафиксировать расширенную таблицу model1 vs model2 по всем радиусам.
\end{itemize}
